\documentclass[../../../include/open-logic-section]{subfiles}

\begin{document}

\olfileid{sth}{ord-arithmetic}{expo}
\olsection{الأس الترتيبي}

ننتقل الآن إلى الأس الترتيبي. وللأسف، لا يوجد تعريف تركيبي
\emph{جميل} للأس الترتيبي.

لا شك أن تعريفات تركيبية صريحة \emph{توجد}. وفيما يأتي أحدها. لتكن
$\text{finfun}(\alpha,\beta)$ مجموعة جميع الدوال
$f \colon \alpha \to \beta$ التي تكون فيها
$\Setabs{\gamma \in \alpha}{f(\gamma) \neq 0}$ مساوية في العدد لعدد
طبيعي ما. وعرّف ترتيبًا جيدًا على~$\text{finfun}(\alpha,\beta)$ بأن
$f \sqsubset g$ إذا وفقط إذا كان~$f \neq g$ و
$f(\gamma_0) < g(\gamma_0)$، حيث
$\gamma_0 = \text{max}\Setabs{\gamma \in \alpha}{f(\gamma) \neq
g(\gamma)}$. ويمكننا عندئذ تعريف~$\ordexpo{\alpha}{\beta}$ بأنه
$\ordtype{\text{finfun}(\beta, \alpha), \sqsubset}$. يستخدم بوتر هذا
التعريف الصريح، ثم يشرح مباشرة:
\begin{quote}
	لا يحدد اختيار هذا الترتيب إلا رغبتنا في الحصول على تعريف للأس
	الترتيبي يحقق الشرط العودي الملائم\ldots، كما أن تصوره أصعب بكثير
	من تصور الجمع المرتب أو الضرب المرتب.
	\citep[p.~199]{Potter2004}
\end{quote}
هذا صحيح تمامًا. فقد شرحنا الجمع بأنه «نسخة من~$\alpha$ تتبعها نسخة
من~$\beta$»، والضرب بأنه «متتالية نمطها~$\beta$ من نسخ~$\alpha$».
لكن ليس لدينا شيء موجز نقوله عن~$\text{finfun}(\gamma, \alpha)$.
ولذلك سنقدم بدلًا من ذلك تعريف الأس الترتيبي \emph{بواسطة} العودية
المتناهية التجاوز فحسب، أي:

\begin{defn}\ollabel{ordexporecursion}
\begin{align*}
	\ordexpo{\alpha}{0} &= 1\\
	\ordexpo{\alpha}{\beta\ordplus 1} &=\ordexpo{\alpha}{\beta} \ordtimes \alpha\\
	\ordexpo{\alpha}{\beta} &= \bigcup_{\delta < \beta}\ordexpo{\alpha}{\delta}& & \text{عندما يكون $\beta$ عددًا ترتيبيًا حديًا}
\end{align*}
\end{defn}

لو كنا نعمل \emph{بوصفنا} منظّري مجموعات، فقد نرغب في استكشاف بعض
خصائص الأس الترتيبي. لكن ليس لدينا الكثير لنضيفه، سوى ملاحظة الحقيقة
غير المفاجئة بأن الأس الترتيبي ليس تبادليًا. وهكذا يكون
$\ordexpo{2}{\omega} =
\bigcup_{\delta < \omega}\ordexpo{2}{\delta} = \omega$، بينما
$\ordexpo{\omega}{2} = \omega \ordtimes \omega$. لكن ينبغي ألا
\emph{نتوقع} أن يكون الأس تبادليًا، لأنه ليس تبادليًا في الأعداد
الطبيعية:~$\ordexpo{2}{3} = 8 < 9 = \ordexpo{3}{2}$.

\begin{prob}
أثبت، باستعمال الاستقراء المتناهي التجاوز، أنه إذا عرّفنا
$\ordexpo{\alpha}{\beta} =
\ordtype{\text{finfun}(\beta, \alpha), \sqsubset}$، فإننا نحصل على
معادلات العودية في
\olref[sth][ord-arithmetic][expo]{ordexporecursion}.
\end{prob}

\end{document}
